% sec_intro.tex -- Introduction. Sources: rigor/phase2_brief.md, rigor/PHASE2_STATUS.md, paper1.
Let $\cA$ be a chiral conformal net on the circle and let $A$, $B$, $C$ be three adjacent
intervals, $B$ the middle one.  The vacuum state $\omega$ restricted to the algebra of
$A\cup B\cup C$ is not a quantum Markov state: the conditional mutual information
$I(A:C\mid B)$ is strictly positive, and, equivalently, the Petz recovery map built from
$\omega|_{\cA(AB)}$ fails to reconstruct $\omega|_{\cA(ABC)}$ from its restriction to
$\cA(AB)$.  The size of that failure is the subject of this paper.

The natural regime is the \emph{zero-collar} one, in which the middle interval is shrunk
away and the recovery defect is governed by a single conformally invariant parameter, the
cross ratio $\zeta$ of the four endpoints.  In Paper~1 \cite{Paper1} we computed the defect
exactly for the chiral free fermion ($c=\tfrac12$ per Majorana component) and found a
quadratic law
\begin{equation}\label{eq:paper1}
 \Phi(\zeta):=-\log \Fid(\omega,\omega_\zeta)=f_2\,\zeta^2+o(\zeta^2),
 \qquad f_2=\frac{c}{12\pi^2},
\end{equation}
where $\Fid$ is the (root) fidelity and $\omega_\zeta$ is the state recovered from the
zero-collar data.  The purpose of the present paper is to show that \eqref{eq:paper1} is
\emph{universal}: the coefficient $f_2$ depends on the net only through its central charge,
and the value $c/(12\pi^2)$ holds for every diffeomorphism covariant conformal net on $S^1$.

\subsection*{Main results}

Our main theorem is stated precisely as \Cref{thm:B} below; informally it reads as follows.

\begin{quote}\itshape
Let $(\cA,U,\Omega)$ be a diffeomorphism covariant conformal net on $S^1$ with central
charge $c$, acting on a separable Hilbert space and satisfying the standard axioms.  Then
the piecewise-M\"obius zero-collar compression of three adjacent intervals is implemented
by unitaries $U_s=e^{isT(g_{\rm tot})}$, and the recovery error obeys
$-\log\Fid(\omega,\omega_s)=\frac{c}{12\pi^{2}}\,\zeta^{2}+o(\zeta^{2})$
in the cross-ratio normalisation $\zeta=s$.
\end{quote}

The proof has three independent components.

\emph{(i) A second-variation formula in type III.}  For a von Neumann algebra $\cM$ with
cyclic separating vector $\Omega$ and a curve of states $\omega_s$ that is differentiable
in the sense of (H3), the logarithm of the fidelity has the second variation
$\frac{s^2}{8}g_B$, where $g_B$ is four times the squared distance of the tangent vector to
the real subspace $\overline{\cM'_{\rm sa}\Omega}$ (\Cref{sec:secvar}).  The lower bound is
Alberti's variational inequality; the upper bound is Uhlmann's inequality applied to
$U_s u'\Omega$ with $u'$ a unitary of the commutant.  The commutant unitaries are exactly
the freedom of extending the compression map beyond the interval, which is why the two
bounds meet.

\emph{(ii) Implementation.}  The compression map $k_s$ is the identity on $A$ and a
M\"obius map on $D$; the two pieces match to first order but not to second, so the glued
map $\tilde k_s$ is $C^{1,1}$ and no better.  Its generating vector field $g_{\rm tot}$
nevertheless has finite $\|\cdot\|_{3/2}$ norm, which is all that the Carpi--Weiner energy
bound \cite{CW} needs; $T(g_{\rm tot})$ is essentially self-adjoint on the finite-energy
domain and the resulting one-parameter group implements $\tilde k_s$ geometrically
(\Cref{sec:impl}).  This holds for \emph{every} diffeomorphism covariant net, with no
restriction on $c$.

\emph{(iii) $c$-linearity and the value of the constant.}  The subspace
$\HT=\overline{\{T(f)\Omega\}}$ generated by the stress tensor reduces the modular data
$(\Delta,J)$ of every interval algebra, because the modular group and the modular
conjugation act on $T$ by M\"obius maps, for which the Schwarzian anomaly vanishes.  The
minimiser of the distance problem therefore lies in $\HT$, where the inner product is $c$
times a net-independent form.  Hence $g_B=c\,\gamma_0$ with a universal $\gamma_0$, and the
free-fermion computation of Paper~1 fixes $\gamma_0=2/(3\pi^2)$ (\Cref{sec:univ}).

\subsection*{What is unconditional and what is not}

The Bures (fidelity) coefficient $f_2=c/(12\pi^2)$ is unconditional for every
diffeomorphism covariant net satisfying the standing assumptions of \Cref{sec:setting}.
The relative-entropy (Kubo--Mori) coefficient is a different matter.  For unitary orbits
generated by an operator \emph{affiliated with the algebra} we prove an exact identity and,
under a modular-analyticity hypothesis, the full second-order expansion; but the generator
$T(g_{\rm tot})$ of the recovery curve is affiliated with neither $\cM$ nor $\cM'$, and the
gauge lemma that would remove this discrepancy is \emph{refuted}: its infimum over admissible gauges is
$(11+5\sqrt5)c/12=11.09\times c/6$, not $c/6$.  We therefore state $s_2=c/12$ as a \emph{conjecture}
(\Cref{sec:km}); what is proved there for a general net is the bracket
$0.034\,c\le\liminf_{s\to0}s^{-2}D\le\limsup_{s\to0}s^{-2}D\le0.924\,c$.
We also record, in \Cref{sec:km}, a route that does \emph{not} work: the Kosaki and
Petz--Donald variational expressions are suprema and yield lower bounds only.

\subsection*{A bosonic check}

Since the argument is structural, a genuinely independent model is a meaningful test.
\Cref{sec:boson} reports a numerical computation of both coefficients in the $U(1)$ current
net (the chiral free boson, $c=1$), a bosonic Gaussian theory sharing no algebraic
structure with the free fermion beyond conformal covariance.  The results
$f_2=0.00844(2)$ and $s_2=0.083(2)$ agree with $1/(12\pi^2)=0.0084434$ and $1/12=0.083333$.
In addition the $c$-linearity of the vacuum overlap is proved analytically at second
order: boson and fermion both give $\frac{s^2}{96\pi^2}\int_0^\infty k^3|\hat g(k)|^2\,dk$.

\subsection*{Notation and organisation}

\Cref{sec:setting} fixes the geometry, the compression map and the hypotheses (H1)--(H3)
with the clause (H2$'$); \Cref{sec:secvar} contains the type-III second-variation lemma;
\Cref{sec:univ} the universality theorem; \Cref{sec:impl} the implementation theorem;
\Cref{sec:km} the Kubo--Mori discussion; \Cref{sec:boson} the bosonic check;
\Cref{sec:outlook} the outlook.  \Cref{app:conv} collects the normalisation conventions,
which matter here because the value of $f_2$ is normalisation dependent, and every constant
in the paper is traceable to that appendix.
